% Unit 1 exam -- AP mirror: 20 MCQ (no calculator) + 2 short FRQ. 48 min.
\documentclass{apchem-exam}

\apcunit{1}
\apcunittitle{Atomic Structure and Properties}
\apcdoctitle{Unit Exam}
\apcsubtitle{Mon Aug 24 \textbullet\ 48 minutes \textbullet\ periodic table provided}

\begin{document}
\apctitleblock

\examsection{I}{Multiple Choice}{20 questions \textbullet\ 30 minutes \textbullet\ 1 point each}{\nocalculator}

% ---- 1.1 moles -------------------------------------------------------------
\begin{mcq}
  What is the mass of \SI{0.50}{\mole} of \ce{H2O}
  ($M = \SI{18}{\gram\per\mole}$)?
  \choices
    \choice{\SI{4.5}{\gram}}
    \correct{\SI{9.0}{\gram}}
    \choice{\SI{18}{\gram}}
    \choice{\SI{36}{\gram}}
\end{mcq}
\why{$0.50 \times 18 = 9.0$.}
\distractor{C}{forgot the 0.50 factor}
\distractor{D}{multiplied by 2 instead of 0.5}

\begin{mcq}
  How many atoms are in \SI{2.0}{\mole} of neon?
  \choices
    \choice{$3.0\times10^{23}$}
    \choice{$6.0\times10^{23}$}
    \correct{$1.2\times10^{24}$}
    \choice{$2.4\times10^{24}$}
\end{mcq}
\why{$2.0 \times 6.0\times10^{23}$.}
\distractor{A}{divided by 2}
\distractor{B}{ignored the 2.0 mol}

\begin{mcq}
  Which sample contains the greatest number of \emph{atoms}?
  \choices
    \correct{\SI{1.0}{\mole} of \ce{NH3}}
    \choice{\SI{2.0}{\mole} of He}
    \choice{\SI{1.0}{\mole} of \ce{CO2}}
    \choice{\SI{3.0}{\mole} of Ar}
\end{mcq}
\why{$1.0 \times 4 = 4.0$ mol of atoms beats 2.0, 3.0, and 3.0.}
\distractor{D}{counted moles of particles, not atoms per particle}

% ---- 1.2 mass spectra ------------------------------------------------------
\begin{mcq}
  An element's mass spectrum shows peaks at 63 u and 65 u with relative
  heights of about 7 to 3. The average atomic mass is closest to
  \choices
    \choice{63.0 u}
    \correct{63.6 u}
    \choice{64.0 u}
    \choice{64.4 u}
\end{mcq}
\why{$0.7(63)+0.3(65) = 63.6$ --- leans toward the abundant isotope.}
\distractor{C}{unweighted average of 63 and 65}

\begin{mcq}
  A mass spectrum of a pure element shows exactly three peaks. This is best
  explained by the element having three
  \choices
    \choice{allotropes}
    \choice{different charges}
    \correct{naturally occurring isotopes}
    \choice{electron shells}
\end{mcq}
\why{One peak per isotope: same Z, different mass.}

\begin{mcq}
  Boron's average atomic mass is 10.8 u, and its isotopes are \ce{^{10}B} and
  \ce{^{11}B}. Which statement must be true?
  \choices
    \choice{The two isotopes are equally abundant.}
    \correct{\ce{^{11}B} is more abundant than \ce{^{10}B}.}
    \choice{\ce{^{10}B} is more abundant than \ce{^{11}B}.}
    \choice{Some boron atoms have a mass of exactly 10.8 u.}
\end{mcq}
\why{10.8 sits much closer to 11: the weighted average leans toward the
abundant isotope.}
\distractor{D}{no individual atom has the average mass}

% ---- 1.3 formulas ----------------------------------------------------------
\begin{mcq}
  The empirical formula of \ce{C6H12O6} is
  \choices
    \choice{\ce{C6H12O6}}
    \choice{\ce{C3H6O3}}
    \correct{\ce{CH2O}}
    \choice{\ce{C2H4O2}}
\end{mcq}
\why{Divide all subscripts by 6 --- smallest whole-number ratio.}

\begin{mcq}
  The percent by mass of oxygen in \ce{CO2} ($M = \SI{44}{\gram\per\mole}$)
  is closest to
  \choices
    \choice{27\%}
    \choice{50\%}
    \choice{67\%}
    \correct{73\%}
\end{mcq}
\why{$32/44 \approx 0.73$.}
\distractor{A}{percent carbon instead}
\distractor{C}{used 2 of 3 atoms rather than mass fractions}

\begin{mcq}
  A compound is 75\% carbon and 25\% hydrogen by mass. Its empirical formula
  is
  \choices
    \choice{\ce{CH}}
    \choice{\ce{CH2}}
    \choice{\ce{CH3}}
    \correct{\ce{CH4}}
\end{mcq}
\why{$75/12 = 6.25$; $25/1 = 25$; ratio $1{:}4$.}
\distractor{C}{ratio of masses 75:25 read as 3:1 C:H}

\begin{mcq}
  Two compounds have the same empirical formula \ce{CH2} but different molar
  masses: \SI{28}{\gram\per\mole} and \SI{56}{\gram\per\mole}. Their
  molecular formulas are
  \choices
    \correct{\ce{C2H4} and \ce{C4H8}}
    \choice{\ce{CH2} and \ce{C2H4}}
    \choice{\ce{C2H4} and \ce{C3H6}}
    \choice{\ce{C2H2} and \ce{C4H4}}
\end{mcq}
\why{$28/14 = 2$ and $56/14 = 4$ empirical units.}

% ---- 1.4 mixtures ----------------------------------------------------------
\begin{mcq}
  A particulate diagram shows some particles that are single filled circles
  and some that are filled--open pairs. The box represents
  \choices
    \choice{a pure element}
    \choice{a pure compound}
    \correct{a mixture of an element and a compound}
    \choice{a single diatomic element}
\end{mcq}
\why{Two distinct particle types = mixture; one type is monatomic element,
the other a two-element compound.}

% ---- 1.5 configurations ----------------------------------------------------
\begin{mcq}
  The ground-state electron configuration of phosphorus is
  \choices
    \correct{$[\ce{Ne}]\,3s^2\,3p^3$}
    \choice{$[\ce{Ne}]\,3s^2\,3p^4$}
    \choice{$[\ce{Ar}]\,3s^2\,3p^3$}
    \choice{$[\ce{Ne}]\,3s^1\,3p^4$}
\end{mcq}
\why{15 electrons; Ar cannot be the core of a lighter element; D violates
Aufbau.}

\begin{mcq}
  Which configuration represents an \emph{excited} state of a neutral
  nitrogen atom?
  \choices
    \choice{$1s^2\,2s^2\,2p^3$}
    \correct{$1s^2\,2s^1\,2p^4$}
    \choice{$1s^2\,2s^2\,2p^4$}
    \choice{$1s^2\,2s^3\,2p^2$}
\end{mcq}
\why{Seven electrons with a 2s vacancy while 2p is occupied. C is oxygen
(8 e$^-$); D is impossible.}

\begin{mcq}
  Which species has the same electron configuration as \ce{Ca^2+}?
  \choices
    \choice{\ce{Na+}}
    \choice{\ce{Ne}}
    \correct{\ce{Cl-}}
    \choice{\ce{Mg^2+}}
\end{mcq}
\why{Both have 18 electrons ($[\ce{Ar}]$).}

% ---- 1.6 PES ---------------------------------------------------------------
\begin{mcq}
  In a photoelectron spectrum, the \emph{area} of a peak indicates the
  \choices
    \choice{binding energy of the electrons}
    \correct{relative number of electrons in that subshell}
    \choice{mass of the isotope}
    \choice{energy of the incoming photons}
\end{mcq}
\why{Position is energy; area is population.}
\distractor{C}{confusing PES with mass spectrometry}

\begin{mcq}
  A complete PES spectrum shows three peaks with relative areas 2 : 2 : 1,
  in order of decreasing binding energy. The element is
  \choices
    \choice{lithium}
    \choice{beryllium}
    \correct{boron}
    \choice{carbon}
\end{mcq}
\why{$1s^2\,2s^2\,2p^1$ --- five electrons.}
\distractor{D}{carbon's last area would be 2}

\begin{mcq}
  The 1s binding energy of Ne is 84 MJ/mol; that of Na is 104 MJ/mol. The
  best explanation is that Na has
  \choices
    \correct{more protons pulling on essentially unshielded 1s electrons}
    \choice{more electron shells than Ne}
    \choice{a larger atomic radius than Ne}
    \choice{fewer valence electrons than Ne}
\end{mcq}
\why{1s sees nearly the full nuclear charge; 11 protons beat 10.}
\distractor{B}{shell count does not lower 1s energy; it is irrelevant to an
unshielded core electron}

% ---- 1.7 trends ------------------------------------------------------------
\begin{mcq}
  Which atom has the largest atomic radius?
  \choices
    \correct{Na}
    \choice{Mg}
    \choice{Al}
    \choice{Si}
\end{mcq}
\why{Same period; smallest $Z_{\text{eff}}$ on the valence shell wins.}

\begin{mcq}
  Successive ionization energies of an element (kJ/mol): 738, 1451, 7733,
  10{,}540. The element is most likely in group
  \choices
    \choice{1}
    \correct{2}
    \choice{13}
    \choice{14}
\end{mcq}
\why{The cliff after $IE_2$ marks two valence electrons (it's Mg).}
\distractor{C}{miscounts the jump position}

% ---- 1.8 ionic -------------------------------------------------------------
\begin{mcq}
  Aluminum and oxygen form an ionic compound with the formula
  \choices
    \choice{\ce{AlO}}
    \choice{\ce{Al2O}}
    \choice{\ce{AlO3}}
    \correct{\ce{Al2O3}}
\end{mcq}
\why{$2(+3) + 3(-2) = 0$.}

\examsection{II}{Free Response}{2 questions \textbullet\ 18 minutes}{\calculatorok}

% ---- FRQ 1: mass spectrum --------------------------------------------------
\frq{short}{4}{CED 1.2 Mass Spectra \textbullet\ SP 2, 4}

Silver's mass spectrum shows two peaks: \SI{106.9}{u} (51.8\%) and
\SI{108.9}{u} (48.2\%).

\begin{parts}
  \item Calculate the average atomic mass of silver. Show your setup.
        \workspace[2.8cm]
        \keyonly{\begin{apcanswer}$0.518(106.9) + 0.482(108.9) =
        \SI{107.9}{u}$\end{apcanswer}}
  \item Explain why the two peaks represent the same element even though the
        masses differ. \answerlines[2]{Both have 47 protons; they differ
        only in neutron count --- isotopes of silver.}
  \item A different sample of silver is found to be enriched in
        \ce{^{109}Ag}. State how the calculated average atomic mass of this
        sample compares to your answer in (a), and why.
        \answerlines[2]{Higher --- greater weight on the 108.9 u peak pulls
        the weighted average upward.}
\end{parts}

\begin{rubric}
  \rpoint{1}{(a) weighted setup using both abundances (as decimals or percents)}
  \rpoint{1}{(a) $107.9 \pm 0.1$ u with units}
  \rpoint{1}{(b) same proton number / isotopes stated explicitly}
  \rpoint{1}{(c) higher average, tied to increased abundance weighting}
\end{rubric}

% ---- FRQ 2: PES + trends ---------------------------------------------------
\frq{short}{4}{CED 1.6, 1.7 PES and Periodic Trends \textbullet\ SP 4, 6}

The complete photoelectron spectrum of magnesium shows peaks at 126, 9.07,
5.31, and 0.74 MJ/mol with relative areas 2 : 2 : 6 : 2.

\begin{parts}
  \item Explain how the spectrum confirms the element is magnesium.
        \answerlines[2]{Areas total $2+2+6+2 = 12$ electrons, matching
        $1s^2 2s^2 2p^6 3s^2$ --- Mg.}
  \item Identify the peak associated with the electrons removed when Mg forms
        \ce{Mg^2+}. \answerlines[1]{The 0.74 MJ/mol peak ($3s^2$) --- the
        least tightly bound.}
  \item Predict whether aluminum's 3s binding energy is greater or less than
        0.74 MJ/mol, and justify with Coulombic reasoning.
        \answerlines[2]{Greater --- Al has one more proton with essentially
        the same shielding of 3s, so $Z_{\text{eff}}$ rises and the electron
        is held more tightly (Al 3s $\approx$ 1.09 MJ/mol).}
\end{parts}

\begin{rubric}
  \rpoint{1}{(a) sums areas to 12 electrons AND links to Mg's configuration}
  \rpoint{1}{(b) selects 0.74 MJ/mol / 3s, with lowest-BE reasoning}
  \rpoint{1}{(c) predicts greater}
  \rpoint{1}{(c) justification uses nuclear charge or $Z_{\text{eff}}$ at
    comparable shielding --- octet or stability language earns nothing}
\end{rubric}

\examtotal

\end{document}
